\section*{Abstract}

Large-scale archetype-based urban building energy modelling relies on building attributes that are often incomplete or missing from open geospatial datasets. This study evaluates whether building type, construction year, and floor count predicted from street view imagery (SVI) can replace reference building attributes from administrative records and reconstructed 3D data in a TABULA-NL building stock enrichment workflow. The experimental dataset comprises 47,150 images linked to 10,086 residential buildings in four Dutch study areas. Three vision configurations were evaluated: a ResNet-50 fine-tuned across all layers, a frozen DINOv2 backbone with supervised prediction heads, and InternVL3-2B without parameter updating. Evaluation spanned building attribute prediction, TABULA-NL archetype assignment, physically weighted differences from reference archetypes, and energy class prediction under a controlled change from reference to predicted attributes.

DINOv2 performed best on six of seven attribute metrics, including a mean absolute error of 9.40 years for construction year. Its joint cell accuracy exceeded both the uniform random expectation and the majority cell baseline. Holding the reconstructed building areas and window-to-wall ratios constant, the aggregate transmission heat transfer coefficient calculated from its predicted archetype cells was approximately 10\% higher than that calculated from reference cells. Replacing reference attributes with DINOv2 predictions had little effect on binary macro F1. However, both the reference attribute condition and the predicted attribute condition had low recall for classes D to G. Direct image classification achieved the highest binary macro F1.

The findings support the conditional use of building attributes predicted from SVI when reference building attributes are unavailable. Suitability depends on the attribute, subsequent task, dataset composition, and geographic relationship between training and application data.

\newpage
\section*{Zusammenfassung}

Großmaßstäbliche urbane Gebäudeenergiemodelle auf Basis von Archetypen benötigen Gebäudeattribute, die in offenen Geodaten häufig unvollständig sind oder fehlen. Diese Arbeit untersucht, ob aus Straßenbildaufnahmen (SVI) vorhergesagte Angaben zu Gebäudetyp, Baujahr und Geschosszahl die Referenzattribute aus Verwaltungsregistern und rekonstruierten dreidimensionalen Daten in einem Verfahren zur Anreicherung des Gebäudebestands mit Archetypen aus TABULA-NL ersetzen können. Der experimentelle Datensatz umfasst 47.150 Bilder, die mit 10.086 Wohngebäuden in vier niederländischen Untersuchungsgebieten verknüpft sind. Bewertet wurden ein in allen Schichten nachtrainiertes ResNet-50, DINOv2 mit eingefrorenen Gewichten und überwachten Vorhersageköpfen sowie InternVL3-2B ohne zusätzliches Training. Die Bewertung umfasste die Vorhersage von Gebäudeattributen, die Zuweisung von Archetypen aus TABULA-NL, physikalisch gewichtete Abweichungen von den Referenzarchetypen und die Klassifikation von Energieklassen unter kontrollierter Ersetzung einzelner Attribute.

DINOv2 erzielte bei sechs von sieben Attributmetriken das beste Ergebnis. Der mittlere absolute Fehler für das Baujahr betrug 9,40 Jahre. Die Genauigkeit der gemeinsamen Archetypzelle übertraf sowohl den Erwartungswert einer gleichverteilten Zufallszuweisung als auch die Baseline der häufigsten Zelle. Bei konstant gehaltenen rekonstruierten Gebäudeflächen und Verhältnissen von Fensterfläche zu Wandfläche war der aus den vorhergesagten Archetypzellen berechnete aggregierte Transmissionskoeffizient etwa 10\% höher als der aus den Referenzzellen berechnete Wert. Die Ersetzung der Referenzattribute durch Vorhersagen von DINOv2 veränderte den binären Makro F1 nur geringfügig. Sowohl die Bedingung mit Referenzattributen als auch die Bedingung mit vorhergesagten Attributen erzielten jedoch eine geringe Sensitivität für die Klassen D bis G. Die direkte Bildklassifikation erreichte beim binären Makro F1 den höchsten Wert.

Die Ergebnisse stützen den bedingten Einsatz von aus Straßenbildaufnahmen abgeleiteten Attributen, wenn keine Referenzdaten verfügbar sind. Die Eignung hängt vom betrachteten Attribut, von der nachgelagerten Aufgabe, von der Zusammensetzung der Stichprobe und von der geografischen Beziehung zwischen den Trainingsdaten und dem Anwendungsgebiet ab.
